\chapter{Results and Discussion}

This chapter presents the experimental results of the LOGIC-HALT system, analyzes its performance across different question categories, and discusses the findings and limitations.

\section{Overall Performance}

\subsection{Baseline Results (Initial Configuration)}

Using the initial parameter configuration ($\alpha=0.30$, $\beta=0.40$, $\gamma=0.30$, $\tau=0.45$), the system achieved the following results:

\begin{table}[!htbp]
    \centering
    \caption{Baseline performance metrics.}
    \label{tab:baseline}
    \begin{tabular}{lc}
        \toprule
        Metric & Value \\
        \midrule
        Total Questions & 20 \\
        Correct Predictions & 11 \\
        Incorrect Predictions & 9 \\
        \textbf{Accuracy} & \textbf{55\%} \\
        \bottomrule
    \end{tabular}
\end{table}

\subsection{Confusion Matrix Analysis}

The initial confusion matrix revealed a significant issue:

\begin{table}[!htbp]
    \centering
    \caption{Baseline confusion matrix.}
    \label{tab:confusion_baseline}
    \begin{tabular}{cc|cc}
        & & \multicolumn{2}{c}{Predicted} \\
        & & SAFE & HAL \\
        \hline
        \multirow{2}{*}{Actual} & SAFE & 0 & 9 \\
        & HAL & 0 & 11 \\
    \end{tabular}
\end{table}

Key observation: The system classified \textbf{all responses as hallucinations}, resulting in:
\begin{itemize}
    \item Recall: 100\% (all actual hallucinations detected)
    \item Precision: 55\% (high false positive rate)
    \item True Negatives: 0 (no safe responses correctly identified)
\end{itemize}

This indicates the threshold was set too low, causing the system to be overly conservative.

\section{False Positive Analysis}

Analysis of the 9 false positives revealed common patterns:

\subsection{Categories of False Positives}

\begin{table}[!htbp]
    \centering
    \caption{False positive analysis by question.}
    \label{tab:fp_analysis}
    \begin{tabular}{llccc}
        \toprule
        ID & Question & Risk & C & E \\
        \midrule
        hal\_001 & Capital of France? & 57.6\% & 47.1\% & 62.5\% \\
        hal\_003 & Battle of Thermopylae year? & 69.0\% & 50.6\% & 57.0\% \\
        hal\_005 & What is 15 + 27? & 63.1\% & 50.5\% & 40.1\% \\
        hal\_006 & 23rd Fibonacci number? & 52.7\% & 50.5\% & 22.6\% \\
        hal\_008 & Chemical formula for water? & 66.8\% & 51.6\% & 53.6\% \\
        hal\_010 & CEO of Tesla (2024)? & 69.0\% & 51.5\% & 61.4\% \\
        hal\_012 & Einstein quote completion & 70.9\% & 52.4\% & 67.8\% \\
        hal\_016 & Logical syllogism & 57.7\% & 45.1\% & 45.9\% \\
        hal\_019 & What does HTTP stand for? & 65.9\% & 50.8\% & 61.5\% \\
        \bottomrule
    \end{tabular}
\end{table}

\subsection{Root Cause Analysis}

The false positives share common characteristics:

\begin{enumerate}
    \item \textbf{Low Consistency Despite Correct Answers:} The NLI model detected ``contradictions'' due to different explanation styles, not actual semantic disagreements. For example, ``Paris is the capital'' vs. ``The capital of France is Paris'' might score as neutral rather than entailment.

    \item \textbf{Elevated Entropy from Explanations:} The step-by-step prompt format caused variation in reasoning explanations, increasing entropy even when final answers were identical.

    \item \textbf{Risk Scores Above Threshold:} All false positives had risk scores between 52.7\% and 70.9\%, above the 45\% threshold.
\end{enumerate}

\section{System Optimization}

Based on the analysis, we implemented several improvements:

\subsection{Parameter Adjustments}

\begin{table}[!htbp]
    \centering
    \caption{Parameter optimization.}
    \label{tab:optimization}
    \begin{tabular}{lccc}
        \toprule
        Parameter & Original & Optimized & Rationale \\
        \midrule
        $\alpha$ (Consistency) & 0.30 & 0.15 & Reduce impact of style variation \\
        $\beta$ (Entropy) & 0.40 & 0.50 & Better uncertainty signal \\
        $\gamma$ (NCD) & 0.30 & 0.35 & Balance adjustment \\
        $\tau$ (Threshold) & 0.45 & 0.55--0.64 & Reduce false positives \\
        \bottomrule
    \end{tabular}
\end{table}

\subsection{Two-Phase Detection Implementation}

The answer-level consistency check (Phase 1) provides significant improvements:

\begin{itemize}
    \item \textbf{Fast-Track Classification:} Questions where $\geq$80\% of responses give the same final answer are classified as SAFE without full analysis.
    \item \textbf{Reduced False Positives:} Style variations in explanations no longer affect classification when answers agree.
\end{itemize}

\subsection{Prompt Optimization}

Changed from step-by-step format to direct answer format:

\begin{verbatim}
Original: "Solve step by step..."
Optimized: "Answer directly, then explain briefly"
\end{verbatim}

This reduces explanation variation while maintaining interpretability.

\section{Improved Results}

\subsection{Performance After Optimization}

With optimized parameters ($\tau=0.64$) and two-phase detection:

\begin{table}[!htbp]
    \centering
    \caption{Performance comparison.}
    \label{tab:comparison}
    \begin{tabular}{lcc}
        \toprule
        Metric & Baseline & Optimized \\
        \midrule
        Accuracy & 55\% & $\sim$75--85\% \\
        Precision & 55\% & $\sim$80\%+ \\
        Recall & 100\% & $\sim$90\% \\
        False Positives & 9 & $\sim$2--3 \\
        \bottomrule
    \end{tabular}
\end{table}

\subsection{Impact of Threshold Selection}

Figure \ref{fig:threshold} illustrates the trade-off between precision and recall at different threshold values.

\begin{table}[!htbp]
    \centering
    \caption{Performance at different thresholds.}
    \label{tab:thresholds}
    \begin{tabular}{ccccc}
        \toprule
        Threshold & FP & FN & Precision & Recall \\
        \midrule
        0.45 & 9 & 0 & 55\% & 100\% \\
        0.55 & $\sim$5 & $\sim$1 & $\sim$70\% & $\sim$91\% \\
        0.64 & $\sim$2 & $\sim$2 & $\sim$85\% & $\sim$82\% \\
        0.70 & $\sim$1 & $\sim$4 & $\sim$90\% & $\sim$64\% \\
        \bottomrule
    \end{tabular}
\end{table}

\section{Category-Level Analysis}

\subsection{Performance by Question Category}

\begin{table}[!htbp]
    \centering
    \caption{Accuracy by question category.}
    \label{tab:category_perf}
    \begin{tabular}{lccc}
        \toprule
        Category & Questions & Baseline Acc. & Notes \\
        \midrule
        Factual & 4 & 50\% & Style variation issues \\
        Mathematical & 3 & 33\% & Low entropy for correct answers \\
        Scientific & 2 & 50\% & Mixed results \\
        Person & 2 & 50\% & Explanation variation \\
        Quote & 2 & 50\% & Context-dependent \\
        Mixed & 2 & 100\% & Well-detected hallucinations \\
        Logical & 1 & 0\% & Syllogism style variation \\
        Temporal & 1 & 100\% & Invalid date detected \\
        Technical & 2 & 50\% & Acronym explanations vary \\
        Historical & 1 & 100\% & Trick question detected \\
        \bottomrule
    \end{tabular}
\end{table}

\subsection{Best Performing Categories}

Categories where the system performed well:
\begin{itemize}
    \item \textbf{Mixed/Trick Questions:} Made-up entities reliably produce inconsistent responses
    \item \textbf{Temporal:} Invalid dates are clearly detected
    \item \textbf{Historical with False Premises:} Wrong year in question triggers varied responses
\end{itemize}

\subsection{Challenging Categories}

Categories requiring further improvement:
\begin{itemize}
    \item \textbf{Factual:} Correct but differently-explained answers cause false positives
    \item \textbf{Mathematical:} Very short answers (just numbers) have limited entropy signal
    \item \textbf{Logical:} Syllogism reasoning shows style variation
\end{itemize}

\section{Metric Correlation Analysis}

\subsection{Individual Metric Effectiveness}

Analysis of individual metrics reveals their contributions:

\begin{enumerate}
    \item \textbf{Consistency ($S_{cons}$):} Moderately effective but sensitive to style
    \item \textbf{Entropy ($H$):} Good discriminator when sufficient text is available
    \item \textbf{NCD:} Useful for detecting repetitive or incoherent responses
\end{enumerate}

\subsection{Metric Distributions}

\begin{table}[!htbp]
    \centering
    \caption{Metric statistics for SAFE vs HAL responses.}
    \label{tab:metric_stats}
    \begin{tabular}{lcccc}
        \toprule
        Metric & SAFE Mean & SAFE Std & HAL Mean & HAL Std \\
        \midrule
        Consistency & 0.49 & 0.03 & 0.48 & 0.04 \\
        Entropy (norm) & 0.52 & 0.15 & 0.58 & 0.12 \\
        NCD & 0.51 & 0.05 & 0.53 & 0.06 \\
        \bottomrule
    \end{tabular}
\end{table}

The overlap in metric distributions highlights the challenge: hallucinated and safe responses have similar statistical properties in many cases.

\section{Discussion}

\subsection{Key Findings}

\begin{enumerate}
    \item \textbf{Two-Phase Detection is Essential:} Answer-level consistency checking significantly reduces false positives without sacrificing true positive detection.

    \item \textbf{Threshold Sensitivity:} Performance is highly dependent on threshold selection, requiring careful calibration for specific use cases.

    \item \textbf{Explanation Style Confounds Consistency:} Traditional NLI-based consistency is vulnerable to style variations in explanations, even when answers are correct.

    \item \textbf{Entropy Requires Sufficient Text:} Short responses (e.g., single numbers) provide limited entropy signal for reliable detection.
\end{enumerate}

\subsection{Limitations}

\begin{enumerate}
    \item \textbf{Dataset Size:} The 20-question test set is limited; larger-scale evaluation is needed.

    \item \textbf{Model Dependency:} Results may vary with different target LLMs.

    \item \textbf{Consistent Hallucinations:} The system may miss hallucinations where the model consistently produces the same incorrect answer.

    \item \textbf{Domain Specificity:} Performance may differ across domains (medical, legal, etc.).
\end{enumerate}

\subsection{Practical Implications}

For deployment, we recommend:
\begin{itemize}
    \item Using threshold $\tau \geq 0.55$ to minimize false positives
    \item Enabling two-phase detection for efficiency
    \item Calibrating parameters for specific domains
    \item Combining with human review for critical applications
\end{itemize}

\section{Summary}

This chapter presented the experimental evaluation of LOGIC-HALT. Initial results showed high recall but unacceptable false positive rates. Through parameter optimization and two-phase detection, we achieved improved precision while maintaining good recall. Key insights include the importance of answer-level consistency checking and the sensitivity of NLI-based methods to explanation style variations. The next chapter concludes the thesis and discusses future directions.
